% =====================================================================
% SUPPLEMENTARY NOTE 1 — why 44 groups of 250 to 500 genes, at 32
% dimensions per token.
%
% SOURCES, and none of this is reconstructed:
%   - AW's own sizing memo, "Reasoning behind model params, group size,
%     and d_token", which is where the two-constraint argument and the
%     d^2 against d^1 asymmetry come from.
%   - Thesis/appendices/C-sizing.tex and D-gene-grouping.tex.
%   - The shipped vocabulary itself, read for every count below:
%     Outputs/supplementary/data/D5_pathway_groups_and_genes.csv.
%
% ⚠ THE THESIS NUMBERS ARE NOT THESE NUMBERS. The thesis build gave 45
% groups over 15,495 genes from a 29,241-gene union. The shipped build
% gives 44 groups over 15,496 genes. Every figure in this note was
% recomputed from the shipped vocabulary; do not copy the thesis table.
% =====================================================================

\section{The gene-group sizing}
\label{note:grouping}

\elaj{} reads a cell as 44 groups of genes rather than as 15{,}496 separate
measurements. Three numbers fix that reading: $D$, how many genes are modelled
at all; $K$, how many groups they fall into, equivalently the median group size
$G \approx D/K$; and $d$, the dimension each gene occupies inside its group.
This note says why those three were set where they were, which is a decision
taken before training and not adjusted since.

\subsection{Optimised variables}

Parameter count is bound by the size of the corpus. With 1{,}577{,}728 cells 
a model of 1{,}883{,}986 parameters is just over one parameter per cell. 
While foundation models have been known to go over two parameters per cell, 
most are closer to one \citep{hao2024scfoundation,cui2024scgpt,theodoris2023geneformer}. Thus, parameter count is the first variable to 
optimise.

The second quantity to optimise here is the dimension of the space that the 
genes are interpreted in. The architecture institutes that each group carries 
its own identifier. An architecture without this rule, $d$ is the maximum 
possible gene space. Therefore, a gene's representation inside group $i$ is 
orthogonal to every representation inside group $j \neq i$. The gene space is
thus block diagonal, and the rank of a block diagonal matrix is the sum of the 
ranks of its blocks. With $K$ groups of $d$ dimensions each, the effective 
gene-space dimension is
\begin{equation}
d_{\mathrm{eff}} = d \cdot K = \frac{d \cdot D}{G}.
\label{eq:deff}
\end{equation}
A gene can be identified along $d_{\mathrm{eff}}$ dimensions making this comparable to 
gene embedding width in other models. gene2vec and the embedding tables of the 
pretrained single-cell models place a gene in 512 dimensions; 
Equation~\ref{eq:deff} is what we have to put beside that figure.

\subsection{The asymmetry that makes the choice}

The cost of a transformer token inside each group is quadratic in the width of 
its tokens \citep{vaswani2017attention}. That implies; parameter count grows with $d^2$
From Equation~\ref{eq:deff}, effective gene space grows with $d^1$. So raising 
$d$ buys representational capacity linearly and cost parameters quadratically.

On the other hand, shrinking $G$ raises $d_{\mathrm{eff}}$ roughly in proportion
to $1/G$, through $K = D/G$ in Equation~\ref{eq:deff}, and costs parameters in
proportion to $1/G$ in the group-encoder term alone.

The reason is that the three parts of the encoder depend on the grouping
differently, and only one of them depends on it at all:
\begin{equation}
P(D, G, d) \;=\; \underbrace{D \cdot d}_{\text{gene embedding}}
\;+\; \underbrace{\frac{D}{G} \cdot E(d)}_{\text{group encoders}}
\;+\; \underbrace{C(d)}_{\text{everything downstream}} .
\label{eq:params}
\end{equation}
Each gene holds one embedding vector of width $d$ wherever it is placed, so the
first term counts every modelled gene once and \textbf{does not depend on the
grouping at all}. The groups are encoded by $K = D/G$ independent
transformers rather than by one shared encoder applied $K$ times, so the second
term is the only one that sees $G$, and $E(d)$ is quadratic in $d$ because the
cost of a transformer layer is. Everything after the group encoders reads a
fixed-width aggregation, so $C(d)$ is independent of $K$ entirely.

Equation~\ref{eq:params} beside Equation~\ref{eq:deff} is the whole argument.
Halving $G$ doubles $d_{\mathrm{eff}}$ and doubles only the middle term of
three. Doubling $d$ also doubles $d_{\mathrm{eff}}$, but it doubles the first
term and quadruples the middle one. Capacity and cost move together in $G$ and
against each other in $d$.

\textbf{So group size, not token width, is the lever that buys representational
capacity cheaply}, and that falls straight out of writing the two quantities
next to each other rather than reasoning about either alone
(Supplementary Fig.~\ref{sfig:params-vs-G} and
Supplementary Fig.~\ref{sfig:params-and-deff}).

\textbf{Two constraints were set in advance.}

\begin{enumerate}
\item \textbf{Fewer than four million parameters}, as a bound set by corpus
size rather than by taste. The corpus is fixed at 1{,}577{,}728 cells and no
more are available, so a model that could memorise it is not a model this
corpus can evaluate.
\item \textbf{An effective gene-space dimension above 512}, because 512 is
where the field's gene embeddings sit. A grouped representation whose effective
dimension fell below that would be buying its inductive bias by giving up
expressiveness, which is a different trade from the one being made.
\end{enumerate}

\subsection{Why the group size is bounded on both sides}

Both bounds are substantive rather than conventional.

\textbf{Above roughly 900 genes per group}, parameter count falls and
$d_{\mathrm{eff}}$ falls faster, so the configuration gets cheaper and less
expressive at the same time. Worse, a larger group is a less meaningful one:
the more genes a group holds, the more dissimilar genes it has had to admit,
and a group that mixes unrelated biology is a group whose summary means less.
The inductive bias the architecture is built on is the claim that a group is a
coherent object, and that claim weakens monotonically in $G$.

\textbf{Below roughly 500 genes per group}, the grouping becomes sharper but
$d$ is forced under 32 to stay inside the parameter bound. That is not a free
trade either. A narrow token with a fine grouping is a stronger bet that the
annotated pathways are the right decomposition of the biology, and for a
disease as under-studied as endometriosis that bet can be counterproductive:
the model is being asked to find a better representation than the annotation
already encodes, and a very fine grouping leaves it little room to.

The usable band is therefore $G$ between about 500 and 900 with $d$ between 32
and 64, and the shipped configuration sits at the sharp end of it.

\subsection{What was built}

Groups are cut from Reactome first and from GO biological process on the
residual, with a minimum group size of 250, a maximum of 500, and genes in
neither hierarchy dropped. The construction is given in full in
Supplementary Methods; Table~\ref{stab:vocabulary} gives the outcome.

\begin{table}[t]
\centering
\caption[The shipped vocabulary]{\textbf{The shipped vocabulary.} Recomputed
from the vocabulary \elaj{} was trained on, not from an earlier build. The
full gene-to-group assignment is Supplementary Data~5.}
\label{stab:vocabulary}
\small
\begin{tabular}{@{}lr@{}}
\toprule
Quantity & Value \\
\midrule
Corpus gene union & 30{,}962 \\
Claimed by Reactome & 10{,}114 genes, 28 groups \\
Claimed by GO from the residual & 5{,}382 genes, 16 groups \\
Dropped, in neither hierarchy & 15{,}466 \\
\midrule
Genes modelled ($D$) & \textbf{15{,}496} \\
Groups ($K$) & \textbf{44} \\
Token dimension ($d$) & \textbf{32} \\
\midrule
Smallest group & 247 genes \\
Median group & 342 genes \\
Mean group & 352 genes \\
Largest group & 493 genes \\
\midrule
Effective gene-space dimension ($d \cdot K$) & \textbf{1{,}408} \\
Model parameters & \textbf{1{,}883{,}986} \\
\bottomrule
\end{tabular}
\end{table}

\textbf{Both constraints are met, and neither is met narrowly.} The model
carries 1{,}883{,}986 parameters against a ceiling of four million, and its
effective gene space is 1{,}408 dimensions against a floor of 512. The second
figure is the one worth stating twice: a representation that reads a cell
through 44 pathway groups is not a smaller gene space than a 512-dimensional
embedding, it is close to three times larger, and the compression the paper
reports is in the \emph{cell} representation at sixteen dimensions rather than
in the gene representation underneath it.

Two features of the table deserve saying plainly rather than leaving to be
noticed. \textbf{Half the corpus union is dropped}: 15{,}466 of 30{,}962 genes
appear in neither hierarchy, and they are excluded on the ground that a gene
with no hierarchical annotation cannot be placed in a group that means
anything. It is a large exclusion and it is a bet, in the same direction as
every other choice here: the model can only find structure among genes somebody
has already annotated. And \textbf{the smallest group holds 247 genes against a
stated floor of 250}, because the floor is applied during the cut and the final
merge can leave one group marginally under it. The discrepancy is three genes
and it is reported rather than rounded away.

\textbf{What the grouping is worth, measured.} The argument above is a
design argument, made before training. The ablation in
the main text is what tests it: replacing the 44 groups with random groups of
identical sizes, parameter count and gene space degraded every biological axis
and disease status most of all, while \emph{improving} held-out reconstruction
to the best figure in the project. The decoder had found a cheaper route and
reconstructed provenance rather than cells.

That is the grouping doing work rather than decorating the input, and it is the
reason this note describes a decision rather than a convention.
